\scanpage{31}
\mspage{30}
\begin{center}-4.11-\end{center}
Now consider the skeleton $\Delta(T)$ in $\Delta(\overline{\mathrm{SL}(3)})$. It can be represented as the set of points on the circle
\[\begin{gathered}r_o+r_1+r_2=0\\r_o^2+r_1^2+r_2^2=6\end{gathered}\]
where the ratios of the $r_i$'s are rational. If $\ell_i$ is the line in $\mathbb P_2$ given by $x_i=0$, and if $P_i$ is the point in $\mathbb P_2$ with coordinates $x_j=\delta_{ij}$, then proceeding around this circle, the parabolic subgroups correspond to flags as indicated in Diagram 5. Furthermore, for a 1-PS $\lambda$ of $T$ as above, put:
\[\Vert\lambda\Vert=\sqrt{(r_o^2+r_1^2+r_2^2)/6}\]
The resulting function $\unsure{\nu|v}^{o(1)}(C,*)$ is indicated in Diagram 6, for some simple curves $C$.

\fbox{\S3.} In this section, we shall consider counter-examples. We would like to show how the reductive groups $\mathrm{SL}(n)$ can operate properly on \uline{quasi-projective} schemes $X$, simple over $\unsure{\mathbb Q|\mathbb C}$, but where there is \uline{no} geometric quotient of $X$ by $\mathrm{SL}(n)$; i.e, by Converse 1.12 and Proposition 0.7, there is no $L$ in $\mathrm{Pic}^G(X)$ such that $X=X^s(L)$. It should be emphasized that in any such case, there is an analytic space which is a quotient of the analytic space defined by $X$,by the analytic group $\mathrm{SL}(n)$: cf. Palais [31]. The moral of these examples is that the concept of stable points is purely algebraic, and not topological.
